\providecommand{\toplevelprefix}{../..}
\documentclass[../../book-main_ko.tex]{subfiles}

\begin{document}

\chapter{비공식적 소개}
\label{ch:intro}

\begin{quote}
``{\em 엔트로피의 지속적인 증가가 우주의 기본 법칙이듯이, 점점 더 높은 구조를 갖추고 엔트로피에 맞서 투쟁하는 것이 삶의 기본 법칙이다.}''

$~$\hfill -- 바츨라프 하벨
 \end{quote}
\vspace{5mm}



\section{지능, 사이버네틱스, 그리고 인공지능}
우리가 사는 세상은 완전히 무작위적이지도 않고 완전히 예측 불가능하지도 않다.\footnote{만약 세상이 완전히 무작위적이라면, 지능적 존재는 어떤 것도 학습하거나 기억할 필요가 없을 것이다.} 대신, 세상은 특정한 질서, 패턴, 법칙을 따르며 이로 인해 대체로 예측 가능하다.\footnote{일부는 결정론적이고, 일부는 확률론적이다.} 생명의 출현과 지속은 바로 이러한 예측 가능성에 의존한다. 환경에서 예측 가능한 것을 학습하고 기억함으로써만 생명은 생존하고 번성할 수 있는데, 건전한 결정과 행동은 신뢰할 수 있는 예측에 달려 있기 때문이다. 세상은 무한히 많은 예측 가능한 현상을 제공하기 때문에, 지능적 존재---동물과 인간---는 점점 더 예민한 감각을 진화시켜왔다: 시각, 청각, 촉각, 미각, 후각. 이러한 감각들은 환경의 규칙성을 인지하기 위해 대용량의 감각 데이터를 수집한다. 따라서 모든 지능적 존재의 근본적인 과제는
\begin{center}
    \textit{대량의 감지된 데이터로부터 예측 가능한 정보를 학습하고 기억하는 것}
\end{center}
이다. 이것이 어떻게 이루어지는지 이해하기 전에, 우리는 세 가지 질문에 답해야 한다:
\begin{itemize}
    \item 예측 가능한 정보를 어떻게 수학적으로 모델링하고 표현할 수 있는가?
    \item 그러한 정보를 데이터로부터 어떻게 효과적이고 효율적으로 계산적으로 학습할 수 있는가?
    \item 이 정보를 미래의 예측과 추론을 지원하기 위해 어떻게 가장 잘 조직해야 하는가?
\end{itemize}
본서는 이러한 질문들에 대한 몇 가지 답을 제공하고자 한다. 이러한 답은 우리가 지능, 특히 지능을 가능하게 하는 계산적 원리와 메커니즘을 더 잘 이해하는 데 도움이 될 것이다. 증거는 모든 형태의 지능---초기 원시 생명체에서 볼 수 있는 저수준 지능부터 현대 과학의 실천이라는 가장 높은 형태의 지능까지---이 공통적인 원리와 메커니즘을 공유한다는 것을 시사한다. 아래에서 이를 자세히 설명한다.

\paragraph{지능의 출현과 진화.}

약 40억 년 전 지구에 생명이 출현하기 위한 필수 조건은 환경이 대체로 예측 가능하다는 것이다. 생명은 환경에 대해 예측 가능한 것을 학습하고, 이 정보를 부호화하며, 생존을 위해 이를 사용할 수 있는 메커니즘을 개발해왔다. 일반적으로 말해서, 우리는 세상에 대한 지식을 학습하는 이러한 능력을 \textit{지능}이라고 부른다. 상당 부분, 생명의 진화는 작동 중인 지능의 메커니즘이다. 생명의 초기 단계에서 지능은 주로 두 가지 유형의 학습 메커니즘을 통해 발전한다: \textit{계통발생적(phylogenetic)} 및 \textit{개체발생적(ontogenetic)}.

\textit{계통발생적 지능}은 종의 진화를 통한 학습을 의미한다. 종은 주로 부모의 DNA 또는 유전자에 부호화된 지식을 기반으로 계승되고 생존한다. 상당 부분, 우리는 DNA를 자연의 사전 훈련된 대형 모델이라고 부를 수 있는데, 그들이 매우 유사한 역할을 하기 때문이다. 계통발생적 지능의 주요 특징은 개체가 제한된 학습 능력을 가진다는 것이다. 학습은 유전자의 무작위 돌연변이를 기반으로 한 ``시행착오'' 메커니즘을 통해 수행되며, 종은 자연 선택---적자생존---에 기반하여 진화한다. 이는 현재 ``강화 학습'' 또는 잠재적으로 ``신경망 아키텍처 탐색''으로 알려진 것의 자연의 구현으로 볼 수 있다. 그러나 이러한 ``시행착오'' 과정은 극도로 느리고, 비용이 많이 들며, 예측 불가능할 수 있다.

% [나머지 섹션들은 번역 예정]
%
% 다음 섹션 구조:
% - 개체발생적 지능 (Ontogenetic Intelligence)
% - 인간 지능의 진화 (Evolution of Human Intelligence)
% - 기계 지능의 기원 - 사이버네틱스 (Origin of Machine Intelligence - Cybernetics)
% - 인공지능의 탄생과 발전 (Birth and Development of Artificial Intelligence)
% - 현대 딥러닝 혁명 (Modern Deep Learning Revolution)
% - 이 책의 구성 (Organization of This Book)
%
% Chapter 1은 975줄로 매우 길어 전체 번역에는 시간이 소요됩니다.
% 서론 부분만 번역하고 나머지는 추후 번역할 예정입니다.

\textit{[나머지 내용은 번역 예정]}

\end{document}
